\documentclass[12pt,a4paper]{report}

\usepackage[utf8]{inputenc}
\usepackage[T1]{fontenc}
\usepackage{lmodern}
\usepackage{setspace}
\usepackage{titlesec}
\usepackage[colorlinks=true, linkcolor=blue, urlcolor=blue, citecolor=blue]{hyperref}



% Rimuove la scritta "Chapter"
\titleformat{\chapter}[hang]
{\normalfont\huge\bfseries}{\thechapter}{1em}{}

\renewcommand{\chaptername}{}

\begin{document}

% =========================
% ABSTRACT
% =========================

\chapter*{Abstract}
\addcontentsline{toc}{chapter}{Abstract}

% Inserire qui l'abstract

\tableofcontents
\newpage

% =========================
% 1 INTRODUCTION
% =========================

\chapter{Introduction}

\section{Context and Motivation}

\section{Research Objectives}

\section{Main Contributions of the Thesis}

\section{Structure of the Document}

% =========================
% 2 THEORETICAL BACKGROUND
% =========================

\chapter{Theoretical Background}

\section{The Evolution of Learning in Artificial Intelligence}

\subsection{Artificial Intelligence Paradigms}

\subsection{Symbolic and Hybrid Approaches in Artificial Intelligence}

\subsection{Reinforcement Learning: Learning Through Experience}

\section{Computer Vision Techniques}

\subsection{Classical Techniques for Video Analysis}

\subsection{Visual Perception in Dynamic Environments}

\subsection{Deep Learning for Feature Extraction from Video Streams}

\section{Knowledge Representation and Reasoning}

\subsection{Symbolic Representations, Semantic Networks, and Interpretable Models}

\section{Reinforcement Learning Foundations}

\subsection{Reinforcement Learning Principles and Deep Q-Networks}

\section{Atari-like Environments as Testbeds}

\subsection{Characteristics and Challenges of 2D Game Environments}

% =========================
% 3 STATE OF THE ART AND CHALLENGES
% =========================

\chapter{State of the Art and Challenges}

\section{Challenges in Symbolic Interpretation}

\subsection{The Illusion of Understanding}

\subsubsection{The Lack of World Models}

\subsection{The Black Box Problem}

\subsubsection{Making Sense of AI: Toward Explainability}

\subsection{Learning Paradigms and Explainable AI}

\subsubsection{Learning Without Understanding}

\subsubsection{Explainable AI: Goals and Approaches}

\subsubsection{Post-hoc Explainability}

\subsubsection{Intrinsic Interpretability}

\subsubsection{Why XAI is Not Enough}

\section{Cognitive Foundations}

\subsection{Human World Modelling and the Origins of Structured Cognition}

\subsection{Core Knowledge Theory: Cognitive Principles and Psychological Foundations}

\subsection{Definition and Relevance for World Understanding}

\subsection{Cognitive Models and Symbolic Representation}

\section{Positioning of the Proposed Framework}

% =========================
% 4 FRAMEWORK DESIGN
% =========================

\chapter{Framework Design - Proposed Framework}

\section{Conceptual Foundations}

\subsection{Evolution of the Idea}

\section{Proposed Implementation}

\subsection{Segmentation}

\subsection{Symbolic System as Knowledge Construction Module}

\subsection{From Symbolic Rules to a Playable Environment}

\subsection{Generic Environment as a Transitional Representation}

\subsection{Deep Q-Network in a Domain-Agnostic Environment}

% =========================
% 5 IMPLEMENTATION
% =========================

\chapter{Implementation}

\section{End-to-End Pipeline}

\section{General Architecture of the Framework}

\subsection{Design Choices}

\subsection{Selected Techniques and Algorithms}

\subsection{Modular Structure of the System}

\subsubsection{Video Acquisition and Preprocessing Module}

\subsubsection{Feature Extraction Module}

\subsubsection{Symbolic Mapping Module - Symbolic Reasoning Engine}

\paragraph{Integration of Core Knowledge for Semantic Understanding}

\paragraph{RL Agent}

\subsection{Optimizations and Efficiency Considerations}

% =========================
% 6 EXPERIMENTAL EVALUATION
% =========================

\chapter{Experimental Evaluation}

\section{Datasets and Testing Scenarios}

\subsection{Evaluation Metrics: Accuracy, Interpretability, Scalability}

\section{Experimental Results}

\subsection{Comparison with Existing Methods}

\subsection{Critical Discussion of the Results}

% =========================
% 7 APPLICATIONS AND CONCLUSION
% =========================

\chapter{Applications and Future Perspectives – Conclusion}

\section{Future Directions}

\subsection{Potential Applications of the Framework}

\subsection{Possible Extensions of the Work}

\subsection{Open Challenges and Future Research Directions}

\section{Summary of Contributions, Impact of the Research, and Final Reflections}

% =========================
% REFERENCES
% =========================

\chapter*{References}
\addcontentsline{toc}{chapter}{References}

% Inserire bibliografia

% =========================
% APPENDICES
% =========================

\appendix

\chapter{Additional Technical Details}

\chapter{Code Snippets or Pseudocode of Relevant Modules}

\chapter{Dataset Specifications or Experimental Settings}

\end{document}